% Part: incompleteness
% Chapter: arithmetization-syntax
% Section: proofs-in-nd

\documentclass[../../../include/open-logic-section]{subfiles}

\begin{document}

% SEGMENT OLP-0287-S01

\olfileid{inc}{art}{pnd}
\olsection{இயல்பான வருவித்தலில் \usetoken{P}{derivation}}

\begin{explain}
!!{derivation}s ஐ எண்கணிதமாக்க, !!{derivation}s ஐ எண்களாகப்
பிரதிநிதித்துவப்படுத்த வேண்டும். ஒவ்வோர் அனுமானமும் ஒன்று அல்லது இரண்டு
சிட்டைகளைத் தாங்கும் !!{formula}s இன் மரங்களே !!{derivation}s என்பதால்,
மீள்வரையறைப் பிரதிநிதித்துவமே மிகத் தெளிவான அணுகுமுறை: ஒரு
!!{derivation} ஐ ஓர் அடுக்காகப் பிரதிநிதித்துவப்படுத்துகிறோம்; அதன் கூறுகள்
கடைசி அனுமானத்தின் முன்கோள்களுக்கு இட்டுச்செல்லும் உடனடி
உட்-!!{derivation}s இன் எண்ணிக்கை, அவற்றின் பிரதிநிதித்துவங்கள்,
இறுதி-!!{formula}, கடைசி அனுமானத்தின் விடுவிப்புச் சிட்டை, மற்றும் கடைசி
அனுமானத்தின் வகையைக் குறிக்கும் ஓர் எண்.
\end{explain}

\begin{defn}
  $\delta$ இயல்பான வருவித்தலில் !!a{derivation} எனில், $\Gn{\delta}$
  பின்வருமாறு தொகுத்தறிதலாக வரையறுக்கப்படுகிறது:
  \begin{enumerate}
  \item
    எடுகோள்~$!A$ மட்டுமே $\delta$ வில் இருந்தால், $\Gn{\delta}$ என்பது
    $\tuple{0, \Gn{!A}, n}$. அது ஓர் !!{undischarged} எடுகோள் எனில்
    எண்~$n$ ஆனது~$0$; இல்லையெனில் அது எண்வடிவச் சிட்டை.
  \item
    பூச்சியம், ஒன்று, இரண்டு அல்லது மூன்று முன்கோள்களைக் கொண்ட ஓர்
    அனுமானத்தில் $\delta$ முடிந்தால், முறையே $\Gn{\delta}$ என்பது
    \begin{align*}
      & \tuple{0, \Gn{!A}, n, k}, \\
      & \tuple{1, \Gn{\delta_1}, \Gn{!A}, n, k}, \\
      & \tuple{2, \Gn{\delta_1}, \Gn{\delta_2}, \Gn{!A}, n, k}, \text{ அல்லது}\\
      & \tuple{3, \Gn{\delta_1}, \Gn{\delta_2}, \Gn{\delta_3}, \Gn{!A}, n,
        k},
    \end{align*}
    இங்கு $\delta_1$, $\delta_2$, $\delta_3$ ஆகியவை $\delta$ வின்
    கடைசி அனுமானத்தின் முன்கோளில் அல்லது முன்கோள்களில் முடியும்
    உட்-!!{derivation}s; $!A$ என்பது $\delta$ வின் கடைசி அனுமானத்தின்
    முடிவுரை; $n$ என்பது கடைசி அனுமானத்தின் விடுவிப்புச் சிட்டை
    (அந்த அனுமானம் எந்த எடுகோளையும் விடுவிக்கவில்லை எனில்~$0$); கடைசி
    அனுமானத்தில் பயன்படுத்தப்பட்ட விதியைப் பொறுத்து பின்வரும் அட்டவணை
    $k$ ஐத் தருகிறது.

    \begin{tabular}{lccccccc}
      \text{விதி:} & \Intro{\land} & \Elim{\land} & \Intro{\lor} & \Elim{\lor} \\
      $k$: & 1 & 2 & 3 & 4  \\[1.5ex]
      \text{விதி:} &  \Intro{\lif} & \Elim{\lif} & \Intro{\lnot} & \Elim{\lnot} \\
      $k$: & 5 & 6 & 7 & 8 \\[1.5ex]
      \text{விதி:}  &  \FalseInt & \FalseCl & \Intro{\lforall} & \Elim{\lforall} \\
      $k$: & 9 & 10 & 11 & 12 \\[1.5ex]
      \text{விதி:} & \Intro{\lexists} & \Elim{\lexists} & \Intro{\eq} & \Elim{\eq} \\
      $k$: & 13 & 14 & 15 & 16
    \end{tabular}
  \end{enumerate}
\end{defn}

\begin{ex}
  மிக எளிய பின்வரும் !!{derivation} ஐக் கருதுக:
  \begin{prooftree}
    \AxiomC{$\Discharge{!A \land !B}{1}$}
    \RightLabel{\Elim{\land}}
    \UnaryInfC{$!A$}
    \DischargeRule{\Intro{\lif}}{1}
    \UnaryInfC{$(!A \land !B) \lif !A$}
  \end{prooftree}
  எடுகோளின் கோடல் எண் $d_0 = \tuple{0,
    \Gn{!A \land !B}, 1}$ ஆக இருக்கும். $\Elim{\land}$ இன் முடிவுரையில்
  முடியும் !!{derivation} இன் கோடல் எண் $d_1 = \tuple{1,
     d_0, \Gn{!A}, 0, 2}$ ஆக இருக்கும் ($\Elim{\land}$ க்கு ஒரு
  முன்கோள் இருப்பதால் $1$; முடிவுரை~$!A$ இன் கோடல் எண்; எந்த எடுகோளும்
  விடுவிக்கப்படாததால் $0$; மேலும் $\Elim{\land}$ ஐக் குறியீடாக்கும் எண்
  $2$). அப்போது முழு !!{derivation} இன் கோடல் எண்
  $\tuple{1, d_1, \Gn{((!A \land !B) \lif
      !A)}, 1, 5}$; அதாவது,
  \[
  \tuple{1, \tuple{1, \tuple{0, \Gn{(!A \land !B)}, 1}, \Gn{!A}, 0, 2},
    \Gn{((!A \land !B) \lif !A)}, 1, 5}.
  \]
\end{ex}

\begin{explain}
!!{derivation}s கான ஒரு பிரதிநிதித்துவத்தைத் தீர்மானித்தபின், அத்தகைய
!!{derivation}s இன் கோடல் எண்களை முதன்மை மீள்வரையறை முறையில் கையாளலாம்
என்றும், அவற்றின் இன்றியமையாப் பண்புகளையும் தொடர்புகளையும் வெளிப்படுத்தலாம்
என்றும் காட்ட வேண்டும். சில செயல்கள் எளியவை: எடுத்துக்காட்டாக,
!!a{derivation} இன் கோடல் எண்~$d$ தரப்பட்டால்,
$\fn{EndFmla}(d) = (d)_{(d)_0+1}$ அதன் இறுதி-!!{formula} வின் கோடல்
எண்ணைத் தருகிறது; $\fn{DischargeLabel}(d) = (d)_{(d)_0 + 2}$ அதன்
விடுவிப்புச் சிட்டையைத் தருகிறது; $\fn{LastRule}(d) = (d)_{(d)_0 + 3}$
கடைசி அனுமானத்தின் வகையைக் குறிக்கும் எண்ணைத் தருகிறது. சில செயல்கள் மிகவும்
கடினமானவை. அவற்றை எவ்வாறு செய்வது என்பதையாவது வரைகோடாகக் காட்டுவோம்.
“$\delta$ என்பது~$\Gamma$ இலிருந்து~$!A$ இன் !!a{derivation}” என்ற
தொடர்பு, $\delta$ மற்றும்~$!A$ இன் கோடல் எண்களின் முதன்மை மீள்வரையறைத்
தொடர்பு என்று காட்டுவதே இலக்கு.
\end{explain}

\begin{prop}
  பின்வரும் தொடர்புகள் முதன்மை மீள்வரையறையானவை:
  \begin{enumerate}
  \item சிட்டை~$n$ உடன் $!A$ ஆனது~$\delta$ வில் ஓர் எடுகோளாக
    இடம்பெறுகிறது.
  \item $\delta$ வில் சிட்டை~$n$ கொண்ட அனைத்து எடுகோள்களும்~$!A$
    வடிவிலுள்ளவை (அதாவது, $\delta$ வில் சிட்டை~$n$ ஐப் பயன்படுத்தி
    எடுகோள்~$!A$ ஐ !!{discharge} செய்யலாம்).
  \end{enumerate}
\end{prop}

\begin{proof}
  !!{formula}s இன் கோடல் எண்களுக்கும் !!{derivation}s இன் கோடல்
  எண்களுக்கும் இடையிலான இவற்றிற்கு ஒத்த தொடர்புகள் முதன்மை
  மீள்வரையறையானவை என்று காட்ட வேண்டும்.
  \begin{enumerate}
  \item கோடல் எண்~$d$ கொண்ட !!{derivation} இல் சிட்டை~$n$ இடப்பட்ட
    ஓர் எடுகோளின் கோடல் எண்~$x$ ஆக இருந்தால் பொருந்தும்
    $\fn{Assum}(x, d, n)$ முதன்மை மீள்வரையறையானது என்று காட்ட விரும்புகிறோம்.
    கோடல் எண்~$\tuple{0, x, n}$ கொண்ட !!{derivation} ஆனது~$d$ இன்
    ஓர் உட்-!!{derivation} ஆக இருந்தால் இது பொருந்தும். !!{derivation}s ஐ
    நாம் குறியீடாக்கும் முறை,
    % READER FALLBACK TA-READER-REF-012
    \oliflabeldef{cmp:rec:tre:sec}{\olref[cmp][rec][tre]{sec}}{மரவுருக் குறியீடாக்கம் பற்றிய முந்தைய பிரிவு}
    இல் அறிமுகப்படுத்திய
    மரங்களின் குறியீடாக்கத்தின் ஒரு சிறப்பு நேர்வு என்பதை நினைவுகூர்க.
    ஆகவே, முதன்மை மீள்வரையறைச் சார்பு $\fn{SubtreeSeq}(d)$ ஆனது~$d$ இன்
    அனைத்து உட்-!!{derivation}s இன் கோடல் எண்களைக் கொண்ட ஒரு தொடரை
    ($d$ ஐ மிஞ்சாத நீளமுடையதை) தருகிறது. எனவே பின்வருமாறு வரையறுக்கலாம்:
    \[
    \fn{Assum}(x, d, n) \defiff \bexists{i<d}{(\fn{SubtreeSeq}(d))_i =
      \tuple{0, x, n}}.
    \]
  \item கோடல் எண்~$d$ கொண்ட !!{derivation} இல் சிட்டை~$n$ கொண்ட
    அனைத்து எடுகோள்களும், கோடல் எண்~$x$ கொண்ட !!{formula} வாக இருந்தால்
    பொருந்தும் $\fn{Discharge}(x, d, n)$ முதன்மை மீள்வரையறையானது என்று
    காட்ட விரும்புகிறோம். ஆனால் இத்தொடர்பு
    $\bforall{y<d}{(\fn{Assum}(y, d, n) \lif y = x)}$ பொருந்தினால்,
    அப்போதும் அப்போது மட்டுமே பொருந்தும்.
  \end{enumerate}
\end{proof}

\begin{prop}
  \ollabel{prop:followsby}
  கோடல் எண்~$d$ கொண்ட !!{derivation}~$\delta$ வின் கடைசி அனுமானம்
  சரியாக இருந்தால், அப்போதும் அப்போது மட்டுமே பொருந்தும்
  $\fn{Correct}(d)$ என்ற பண்பு முதன்மை மீள்வரையறையானது.
\end{prop}

\begin{proof}
  ஒவ்வோர் அனுமான விதி~$R$ க்கும் $\fn{FollowsBy}_R(d)$ என்ற தொடர்பு
  முதன்மை மீள்வரையறையானது என்று இங்கு காட்ட வேண்டும். $d$ ஆனது
  !!{derivation}~$\delta$ இன் கோடல் எண்ணாகவும், $\delta$ வின்
  இறுதி-!!{formula} ஆனது அதன் உடனடி உட்-!!{derivation}s இலிருந்து~$R$ ஐச்
  சரியாகப் பயன்படுத்திப் பின்தொடர்ந்ததாகவும் இருந்தால், அப்போதும் அப்போது
  மட்டுமே $\fn{FollowsBy}_R(d)$ பொருந்தும்.

  \Intro{\land} விதி ஓர் எளிய நேர்வு. சரியான $\Intro{\land}$
  அனுமானத்தில் $\delta$ முடிந்தால், அது பின்வருமாறு தோன்றும்:
  \begin{prooftree}
    \AxiomC{}
    \RightLabel{$\delta_1$}
    \DeduceC{$!A$}

    \AxiomC{}
    \RightLabel{$\delta_2$}
    \DeduceC{$!B$}

    \RightLabel{\Intro\land}
    \BinaryInfC{$!A \land !B$}
  \end{prooftree}
  அப்போது $\delta$ வின் கோடல் எண்~$d$ என்பது $\tuple{2, d_1, d_2,
    \Gn{(!A \land !B)}, 0, k}$; இங்கு $\fn{EndFmla}(d_1) = \Gn{!A}$,
  $\fn{EndFmla}(d_2) = \Gn{!B}$, $n=0$, மற்றும் $k=1$. எனவே
  $\fn{FollowsBy}_{\Intro\land}(d)$ ஐப் பின்வருமாறு வரையறுக்கலாம்:
  \begin{multline*}
    (d)_0 = 2 \land \fn{DischargeLabel}(d) = 0 \land \fn{LastRule}(d) = 1 \land {}\\
    \fn{EndFmla}(d) = {}\\
    \Gn{(} \concat \fn{EndFmla}((d)_1) \concat \Gn{\land}
    \concat \fn{EndFmla}((d)_2) \concat \Gn{)}.
  \end{multline*}

  மற்றோர் எளிய எடுத்துக்காட்டு $\Intro\eq$ விதி. எடுகோள்களைப் போலவே,
  இதற்கு முன்கோள்கள் இல்லை; ஆகவே $(d)_0 = 0$. இதற்கு விடுவிப்புச்
  சிட்டையும் இல்லை; அதாவது, $n=0$. எனினும், ஏதாவது மூடிய சொல்~$t$ க்கு
  $!A$ ஆனது $\eq[t][t]$ வடிவில் இருக்க வேண்டும். இங்கு ஒரு முதன்மை
  மீள்வரையறை வரையறை:
  \begin{multline*}
    (d)_0 = 0 \land \fn{DischargeLabel}(d) = 0 \land {}\\
    \bexists{t<d}{(\fn{ClTerm}(t) \land
      \fn{EndFmla}(d) = {}\\
      \Gn{{\eq}(} \concat t \concat \Gn{,} \concat t \concat \Gn{)})}.
  \end{multline*}

  சற்றுச் சிக்கலான எடுத்துக்காட்டாக,
  $\fn{FollowsBy}_{\Intro{\lif}}(d)$ ஐக் கருதுக. $\delta$ வின்
  இறுதி-!!{formula} $(!A \lif !B)$ வடிவிலும், $\delta_1$ இன்
  இறுதி-!!{formula}~$!B$ ஆகவும், சிட்டை~$n$ உடைய $\delta$ வின் ஒவ்வோர்
  எடுகோளும்~$!A$ வடிவிலும் இருந்தால், அப்போதும் அப்போது மட்டுமே இது
  பொருந்தும். இதைப் பின்வருமாறு முதன்மை மீள்வரையறை முறையில் வெளிப்படுத்தலாம்:
  % SOURCE CORRECTION TA-INC-008: discharge label 0 denotes vacuous discharge.
  \begin{multline*}
    (d)_0 = 1 \land {}\\
    \bexists{a<d}{((\fn{DischargeLabel}(d) = 0 \lor
      \fn{Discharge}(a, (d)_1, \fn{DischargeLabel}(d))) \land {}}\\
      \fn{EndFmla}(d) = (\Gn{(} \concat a \concat \Gn{\lif}
      \concat \fn{EndFmla}((d)_1) \concat \Gn{)}))
  \end{multline*}
  ($a$ ஐ~$!A$ இன் கோடல் எண்ணாகக் கருதுக).

  மற்றோர் எடுத்துக்காட்டாக, \Intro{\lexists} ஐக் கருதுக.
  !!a{formula}~$!A$, ஒரு மூடிய சொல்~$t$, மற்றும் !!a{variable}~$x$
  இருந்து, $\Subst{!A}{t}{x}$ ஆனது !!{derivation}~$\delta_1$ இன்
  இறுதி-!!{formula} ஆகவும், $\lexists[x][!A]$ கடைசி அனுமானத்தின்
  முடிவுரையாகவும் இருந்தால், அப்போதும் அப்போது மட்டுமே $\delta$ வின்
  கடைசி அனுமானம் சரியானது. எனவே,
  $\fn{FollowsBy}_{\Intro{\lexists}}(d)$ பொருந்துவது, அப்போதும் அப்போது
  மட்டுமே,
  \begin{multline*}
    (d)_0 = 1 \land \fn{DischargeLabel}(d) = 0 \land {} \\
    \bexists{a < d}{\bexists{x<d}{\bexists{t<d}{
          (\fn{ClTerm}(t) \land \fn{Var}(x)  \land {}}}}\\
    \fn{Subst}(a,t,x) = \fn{EndFmla}((d)_1) \land
    \fn{EndFmla}(d) = (\Gn{\lexists} \concat x \concat a)).
  \end{multline*}

  அடுத்து, $\fn{Correct}(d)$ ஐப் பின்வருமாறு வரையறுக்கிறோம்:
  % SOURCE CORRECTION TA-INC-006: group every rule/assumption branch under Sent.
  \begin{multline*}
    \fn{Sent}(\fn{EndFmla}(d)) \land {}\\
    [ (\fn{LastRule}(d) = 1 \land
    \fn{FollowsBy}_{\Intro\land}(d)) \lor \dots \lor {}\\
    (\fn{LastRule}(d) = 16 \land \fn{FollowsBy}_{\Elim\eq}(d)) \lor {}\\
    \bexists{n<d}{\bexists{x<d}{(d = \tuple{0, x, n})}}].
  \end{multline*}
  முதல் வரி, $d$ இன் இறுதி-!!{formula} ஒரு வாக்கியம் என்பதை
  உறுதிசெய்கிறது. $d$ ஓர் எடுகோள் மட்டுமே ஆகும் நேர்வைக் கடைசி வரி
  கையாளுகிறது.
\end{proof}

\begin{prob}
  \olref[inc][art][pnd]{prop:followsby} இல் உள்ளவாறு பின்வரும் பண்புகளை
  வரையறுக்க:
  \begin{enumerate}
  \item $\fn{FollowsBy}_{\Elim{\lif}}(d)$,
  \item $\fn{FollowsBy}_{\Elim{\eq}}(d)$,
  \item $\fn{FollowsBy}_{\Elim{\lor}}(d)$,
  \item $\fn{FollowsBy}_{\Intro{\lforall}}(d)$.
  \end{enumerate}
  கடைசிப் பண்புக்கு, கோடல் எண்~$d$ கொண்ட !!{derivation} இன் கடைசி
  அனுமானம் தனிமாறி நிபந்தனையை நிறைவுசெய்கிறதா என்பதையும் முதன்மை
  மீள்வரையறை முறையில் சோதிக்கலாம் என்று காட்ட வேண்டும்; அதாவது,
  $\Intro{\lforall}$ அனுமானத்தின் தனிமாறி~$a$, $d$ இன்
  இறுதி-!!{formula} விலோ $d$ இன் திறந்த எடுகோளிலோ இடம்பெறக்கூடாது.
  இதற்கு \olref[inc][art][pnd]{prop:openassum} இலுள்ள முதன்மை
  மீள்வரையறைப் பயனிலை $\fn{OpenAssum}$ ஐப் பயன்படுத்தலாம்.
\end{prob}

\begin{prop}
  \ollabel{prop:deriv}
  $d$ ஆனது சரியான !!{derivation}~$\delta$ இன் கோடல் எண்ணாக இருந்தால்
  பொருந்தும் $\fn{Deriv}(d)$ என்ற தொடர்பு முதன்மை மீள்வரையறையானது.
\end{prop}

\begin{proof}
  !!^a{derivation}~$\delta$ வின் ஒவ்வோர் அனுமானமும் ஒரு விதியின் சரியான
  பயன்பாடாக, அதாவது அதன் ஒவ்வோர் உட்-!!{derivation} உம் ஒரு சரியான
  அனுமானத்தில் முடிந்தால், $\delta$ சரியானது. ஆகவே,
  $\fn{Deriv}(d)$ பொருந்தினால், அப்போதும் அப்போது மட்டுமே,
  \[
  \bforall{i<\len{\fn{SubtreeSeq}(d)}}{\fn{Correct}((\fn{SubtreeSeq}(d))_i)}
  \]
\end{proof}

\begin{prop}
  \ollabel{prop:openassum} கோடல் எண்~$d$ கொண்ட
  !!{derivation}~$\delta$ வின் !!a{undischarged} எடுகோள்~$!A$ இன்
  கோடல் எண்~$z$ ஆக இருந்தால் பொருந்தும் $\fn{OpenAssum}(z, d)$ என்ற
  தொடர்பு முதன்மை மீள்வரையறையானது.
\end{prop}

\begin{proof}
  ஓர் எடுகோளின் நிகழ்வு, சிட்டை~$n$ உடன் $\delta$ வின் ஓர்
  உட்-!!{derivation} இல் இடம்பெற்று, சிட்டை~$n$ ஐ விடுவிப்புச் சிட்டையாகக்
  கொண்ட ஒரு விதியில் அந்த உட்-!!{derivation} முடிந்தால்,
  !!{discharged} ஆகும். ஆகவே, அவ்வெடுகோளின் நிகழ்வுகளில் குறைந்தது ஒன்று
  $\delta$ வில் !!{discharged} ஆகவில்லை எனில், $!A$ ஆனது~$\delta$ வின்
  !!a{undischarged} எடுகோள். இங்கு கவனம் தேவை: $\delta$ ஆனது~$!A$ இன்
  !!{discharged} நிகழ்வுகளையும் !!{undischarged} நிகழ்வுகளையும்
  ஒருங்கே கொண்டிருக்கலாம்.

  $\delta_0 = \delta$ ஆகவும், $\delta_k$ ஆனது ஏதாவது~$n$ க்கான
  எடுகோள் $\Discharge{!A}{n}$ ஆகவும், $\delta_{i+1}$ ஆனது
  $\delta_i$ இன் உடனடி உட்-!!{derivation} ஆகவும் இருக்கும் ஒரு தொடர்
  $\delta_0$, \dots, $\delta_k$ ஐக் கருதுக. எந்த~$\delta_i$ உம்
  சிட்டை~$n$ ஐ விடுவிப்புச் சிட்டையாகக் கொண்ட ஓர் அனுமானத்தில் முடியாத
  இத்தகைய தொடர் இருந்தால், $!A$ ஆனது~$\delta$ வின்
  !!a{undischarged} எடுகோள்.

  முதன்மை மீள்வரையறைச் சார்பு $\fn{SubtreeSeq}(d)$ ஆனது~$\delta$ வின்
  அனைத்து உட்-!!{derivation}s இன் கோடல் எண்களைக் கொண்ட ஒரு தொடரைத்
  தருகிறது. $\delta$ வின் உட்-!!{derivation}s இன் கோடல் எண்களால் ஆன எந்தத்
  தொடரும் அதன் ஒரு உட்தொடர். உட்தொடராக இருப்பது முதன்மை மீள்வரையறைத்
  தொடர்பு: $\fn{Subseq}(s, s')$ பொருந்துவது, அப்போதும் அப்போது மட்டுமே,
  $\bforall{i<\len{s}}{\lexists[j<\len{s'}][(s)_i = (s')_j]}$.
  உடனடி உட்-!!{derivation} ஆக இருப்பதும் அவ்வாறே:
  % SOURCE CORRECTION TA-INC-007: tuple slots 1 through arity hold immediate subderivations.
  $\fn{Subderiv}(d, d')$ பொருந்துவது, அப்போதும் அப்போது மட்டுமே,
  $\bexists{j<(d')_0}{d = (d')_{j+1}}$. எனவே
  $\fn{OpenAssum}(z, d)$ ஐப் பின்வருமாறு வரையறுக்கலாம்:
  % SOURCE CORRECTION TA-INC-009: label 0 explicitly marks an undischarged assumption.
  \begin{multline*}
    \bexists{s<\fn{SubtreeSeq}(d)}{(\fn{Subseq}(s, \fn{SubtreeSeq}(d))
      \land (s)_0 = d \land {}} \\
    \bexists{n<d}{((s)_{\len{s} \tsub 1} = \tuple{0, z, n} \land {}}\\
      \bforall{i<(\len{s} \tsub 1)}{(\fn{Subderiv}((s)_{i+1}, (s)_i) \land {}}\\
      (n = 0 \lor \fn{DischargeLabel}((s)_i) \neq n)))).
  \end{multline*}
\end{proof}

\begin{prop}
  \ollabel{prop:prf-prim-rec}
  $\Gamma$ என்பது !!{sentence}s இன் முதன்மை மீள்வரையறைக் கணம் என்க.
  அப்போது, “$x$ என்பது~$\Gamma$ விலுள்ள !!{undischarged} எடுகோள்களிலிருந்து
  $!A$ க்கான !!a{derivation}~$\delta$ வின் குறியீடாகவும், $y$ என்பது~$!A$
  இன் கோடல் எண்ணாகவும் உள்ளது” என்பதை வெளிப்படுத்தும்
  $\Prf[\Gamma](x, y)$ என்ற தொடர்பு முதன்மை மீள்வரையறையானது.
\end{prop}

\begin{proof}
  “$y \in \Gamma$” என்பது $R_\Gamma(y)$ என்ற முதன்மை மீள்வரையறைப்
  பயனிலையால் தரப்படுகிறது என்க. $y$ என்பது ஒரு வாக்கியம்~$!A$ இன்
  கோடல் எண்ணாகவும், $x$ என்பது இறுதி-!!{formula}~$!A$ ஐயும்
  $\Gamma$ விலுள்ள அனைத்து !!{undischarged} எடுகோள்களையும் கொண்ட இயல்பான
  வருவித்தல் !!{derivation} இன் குறியீடாகவும் இருந்தால், அப்போதும் அப்போது
  மட்டுமே பொருந்தும் $\Prf[\Gamma](x, y)$ முதன்மை மீள்வரையறையானது என்று
  காட்ட வேண்டும்.

  \olref{prop:deriv} இன்படி, $x$ ஆனது இயல்பான வருவித்தலில் சரியான
  !!{derivation}~$\delta$ இன் கோடல் எண்ணாக இருந்தால், அப்போதும் அப்போது
  மட்டுமே பொருந்தும் $\fn{Deriv}(x)$ என்ற பண்பு முதன்மை
  மீள்வரையறையானது. ஆகவே, $\Prf[\Gamma](x, y)$ ஐப் பின்வருமாறு
  வரையறுக்கலாம்:
  \begin{align*}
    \Prf[\Gamma](x, y) \defiff {}
    & \fn{Deriv}(x) \land \fn{EndFmla}(x) = y \land {} \\
    & \bforall{z < x}{(\fn{OpenAssum}(z, x) \lif R_\Gamma(z))}.
  \end{align*}
\end{proof}

\end{document}
